% ============================================================================
% CAPÍTULO 6 — TypeScript\index{TypeScript}
% Texto completo (rodada editorial 2)
% ============================================================================
\chapter{TypeScript: Tipos para Código que Escala}
\label{cap:typescript}

\begin{caixaconceito}
\textbf{TypeScript} é um superconjunto tipado de JavaScript: todo JavaScript
válido é TypeScript válido. Ele adiciona \emph{tipos estáticos} que são
removidos na compilação, deixando JavaScript puro para o navegador — e, desde
agosto de 2025, é a linguagem mais usada do GitHub \cite{octoverse2025}.
\end{caixaconceito}

Ao final deste capítulo, você será capaz de:
\begin{objetivos}
  \item Explicar por que tipos estáticos transformam a forma de programar;
  \item Configurar um projeto TypeScript com \texttt{tsconfig.json} em \emph{strict mode};
  \item Tipar variáveis, funções, objetos e arrays com inferência e anotações;
  \item Modelar domínios com \texttt{interface}, \texttt{type}, unions e literals;
  \item Escrever \emph{generics} reutilizáveis e fazer \emph{narrowing} seguro;
  \item Aplicar tipos a código de DOM e a respostas de APIs;
  \item Entregar o projeto do capítulo: uma biblioteca de validação tipada.
\end{objetivos}

\section{Por que tipos? O custo do JavaScript sem rede de segurança}

JavaScript é uma linguagem \emph{dinâmica}: uma função aceita qualquer valor,
e o erro só aparece quando o dado errado chega — muitas vezes em produção, com
um usuário real na frente. O TypeScript move esse erro para o \emph{tempo de
compilação}: o editor marca o problema em vermelho \emph{antes} de você rodar.

Pense em tipos como a \emph{planta do prédio} antes da construção: você
descobre que a porta não cabe na parede no papel, e não depois de assentar os
tijolos. Em um projeto com milhares de arquivos — e com assistentes de IA
gerando código a cada tecla —, tipos são a diferença entre refatorar com
confiança e refatorar com medo.

\begin{caixadica}
O editor (VS Code) usa os tipos para dar \emph{autocomplete}, renomear
símbolos com segurança e mostrar documentação na hora. Quem programa com
TypeScript ``anda de bicicleta com rodinhas'' — mas as rodinhas são a planta
do prédio: elas impedem a queda e documentam o projeto.
\end{caixadica}

\section{Instalação e configuração}

Crie uma pasta para o projeto e rode:

\begin{lstlisting}[style=livro, language=Bash,
  caption={Iniciando um projeto TypeScript.}, label={list:ts-init}]
mkdir tipos-na-pratica && cd tipos-na-pratica
npm init -y
npm install -D typescript @types/node
npx tsc --init              # gera o tsconfig.json
\end{lstlisting}

O comando \texttt{tsc --init} gera um \texttt{tsconfig.json} com dezenas de
opções comentadas. O coração da configuração é o \emph{strict mode}:

\begin{lstlisting}[style=livro, language=JSON,
  caption={Opções essenciais do tsconfig.json.}, label={list:ts-config}]
{
  "compilerOptions": {
    "strict": true,
    "target": "ES2023",
    "module": "ESNext",
    "moduleResolution": "bundler",
    "noEmit": true,
    "verbatimModuleSyntax": true,
    "skipLibCheck": true
  },
  "include": ["src"]
}
\end{lstlisting}

\begin{caixadica}
Sempre mantenha \texttt{"strict": true}. É ele que liga
\texttt{strictNullChecks} — a opção que elimina a classe inteira de bugs de
\texttt{null} e \texttt{undefined} que assombra o JavaScript. Em projetos
legados sem strict\index{strict}, a migração começa ligando essa opção e corrigindo os
erros um a um.
\end{caixadica}

\begin{caixaconceito}
\textbf{Como o TypeScript executa?} Ele \emph{compila} para JavaScript puro: os
tipos são apagados (``type erasure'') e o resultado roda em qualquer runtime.
Por isso \texttt{noEmit: true} é comum em projetos que delegam a execução a
ferramentas (Vite, Next.js, Node com type stripping): o \texttt{tsc} serve
apenas para \emph{checar} tipos.
\end{caixaconceito}

\section{Tipos básicos e inferência}

O TypeScript \emph{infere} tipos quando você não anota: ele lê o valor inicial
e deduz o tipo. A inferência é a base de toda a experiência — você escreve
JavaScript e ganha tipos ``de graça'':

\begin{lstlisting}[style=livro, language=TypeScript,
  caption={Inferência e anotações explícitas.}, label={list:ts-basico}]
// Inferência: o TypeScript descobre o tipo sozinho
let nome = "Ana";          // string
let idade = 28;            // number
let ativo = true;          // boolean

// Erro de inferência: tentar mudar o tipo quebra na hora
// nome = 42;              // Erro: type 'number' is not assignable to 'string'

// Anotação explícita: útil em parâmetros, retornos e fronteiras
function calcularTotal(valor: number, desconto: number): number {
  return valor * (1 - desconto);
}

// Tipos primitivos completos
const texto: string = "olá";
const numero: number = 19.9;
const ligado: boolean = true;
const nulo: null = null;
let indefinido: undefined;
const simbolo: symbol = Symbol("id");
const grande: bigint = 9007199254740993n;
\end{lstlisting}

\begin{caixadica}
Regra prática: \emph{anote as fronteiras, infira o interior}. Anote
parâmetros de função, retornos de APIs e variáveis sem valor inicial; deixe a
inferência cuidar do resto. Anotar tudo é ruído; anotar nada é imprudência.
\end{caixadica}

\section{Arrays, objetos e tuplas}

\begin{lstlisting}[style=livro, language=TypeScript,
  caption={Arrays, objetos e tuplas tipados.}, label={list:ts-colecoes}]
// Arrays: duas sintaxes equivalentes
const numeros: number[] = [1, 2, 3];
const nomes: Array<string> = ["Ana", "Bia"];

// Objetos: shape inline
const usuario: { nome: string; idade: number } = {
  nome: "Ana",
  idade: 28,
};

// Tuplas: array com posições e tipos fixos
const parCoordenadas: [number, number] = [-23.55, -46.63];

// readonly: imutabilidade no nível do tipo
const config: readonly string[] = ["produção"];
// config.push("dev"); // Erro: property 'push' does not exist
\end{lstlisting}

\begin{caixaconceito}
\textbf{Structural typing}: o TypeScript compara \emph{estruturas}, não nomes.
Se dois objetos têm o mesmo shape, eles são compatíveis — mesmo que venham de
interfaces diferentes. Isso é o que permite tipar APIs de terceiros sem
herança, e é diferente do \emph{nominal typing} de linguagens como Java.
\end{caixaconceito}

\section{\texttt{interface} vs \texttt{type}: modelando o domínio}

A decisão mais comum do dia a dia: quando usar \texttt{interface} e quando
usar \texttt{type}?

\begin{lstlisting}[style=livro, language=TypeScript,
  caption={interface, type e composição.}, label={list:ts-interface}]
// interface: contrato extensível (herança, implementação)
interface Usuario {
  id: number;
  nome: string;
  email: string;
  ativo?: boolean;          // propriedade opcional
  readonly criadoEm: Date;  // propriedade somente leitura
}

// type: composição e unions
type Papel = "admin" | "editor" | "leitor";
type Id = string | number;

interface Membro extends Usuario {
  papel: Papel;
  permissoes: string[];
}

// Intersecção com type
type MembroCompleto = Usuario & { papel: Papel };

const ana: Membro = {
  id: 1,
  nome: "Ana Souza",
  email: "ana@exemplo.com",
  papel: "editor",
  permissoes: ["ler", "escrever"],
  criadoEm: new Date(),
};
\end{lstlisting}

\begin{caixaconceito}
\texttt{interface} é melhor para \emph{contratos que se estendem}
(\texttt{extends}, \texttt{implements}, \emph{declaration merging}); \texttt{type}
é melhor para \emph{composições} (unions, intersecções, tipos utilitários,
tuplas nomeadas). Ambos coexistem; times modernos escolhem um padrão e seguem
a convenção — o importante é a consistência.
\end{caixaconceito}

\section{Unions, literals e narrowing}

\begin{lstlisting}[style=livro, language=TypeScript,
  caption={Unions, literals e narrowing com segurança.}, label={list:ts-union}]
// Union de literais: o valor só pode ser um dos listados
type StatusPedido = "pendente" | "pago" | "enviado" | "entregue";

function descrever(status: StatusPedido): string {
  // narrowing: o TypeScript restringe o tipo em cada branch
  switch (status) {
    case "pendente":
      return "Aguardando pagamento";
    case "pago":
      return "Pagamento confirmado";
    case "enviado":
      return "A caminho";
    case "entregue":
      return "Entregue";
  }
}

// Union de tipos: string OU number
function formatarId(id: string | number): string {
  if (typeof id === "string") {
    return id.toUpperCase();      // aqui id é string
  }
  return id.toFixed(0);           // aqui id é number
}

// Narrowing com in / instanceof
interface Carro { rodas: 4; motor: string }
interface Moto { rodas: 2 }

function ehCarro(veiculo: Carro | Moto): veiculo is Carro {
  return "motor" in veiculo;
}
\end{lstlisting}

\begin{caixaconceito}
\textbf{Narrowing} é o processo de restringir o tipo de uma variável dentro de
um bloco: \texttt{typeof}, \texttt{instanceof}, \texttt{in}, \texttt{switch}
e comparações com literals. O TypeScript acompanha o fluxo do código e sabe —
em cada linha — qual tipo é possível. É assim que você transforma ``qualquer
coisa'' em ``exatamente o que preciso''.
\end{caixaconceito}

\section{Discriminated unions: o padrão de modelagem mais poderoso}

Quando objetos têm uma propriedade \emph{discriminante} (geralmente
\texttt{tipo}), o narrowing\index{narrowing} funciona por \texttt{switch} completo:

\begin{lstlisting}[style=livro, language=TypeScript,
  caption={Discriminated union para resultados de operações.},
  label={list:ts-discriminated}]
type Resposta<T> =
  | { tipo: "sucesso"; dados: T; tempo: number }
  | { tipo: "erro"; mensagem: string; codigo: number };

function tratar<T>(resposta: Resposta<T>): string {
  switch (resposta.tipo) {
    case "sucesso":
      // aqui o TypeScript SABE que existe resposta.dados
      return `OK em ${resposta.tempo}ms com ${JSON.stringify(resposta.dados)}`;
    case "erro":
      return `Falha ${resposta.codigo}: ${resposta.mensagem}`;
  }
}
\end{lstlisting}

\begin{caixadica}
Discriminated unions são a forma padrão de modelar \emph{resultados de
operações}, \emph{estados de UI} (carregando/sucesso/erro) e \emph{mensagens
de sistemas distribuídos}. Você os verá em bibliotecas profissionais e em
entrevistas — domine este padrão.
\end{caixadica}

\section{Generics: código que se adapta ao tipo}

Generics\index{generics} permitem escrever uma função que funciona \emph{para qualquer tipo},
preservando o tipo de entrada na saída:

\begin{lstlisting}[style=livro, language=TypeScript,
  caption={Generics: um único código, todos os tipos.}, label={list:ts-generic}]
function primeiroItem<T>(lista: T[]): T | undefined {
  return lista[0];
}

const numero = primeiroItem([10, 20, 30]);     // number | undefined
const nome = primeiroItem(["ana", "bia"]);     // string | undefined

// Dois parâmetros de tipo
function chavear<Chave, Valor>(chave: Chave, valor: Valor): Map<Chave, Valor> {
  const mapa = new Map<Chave, Valor>();
  mapa.set(chave, valor);
  return mapa;
}

// Constraints: T precisa ter certa forma
function tamanhoDaChave<T extends { id: number }>(item: T): number {
  return String(item.id).length;
}

// Arrow function genérica
const identidade = <T,>(x: T): T => x;
\end{lstlisting}

\begin{caixaconceito}
Um \textbf{generic} (\texttt{<T>}) é um ``buraco de tipo'' que o compilador
preenche em cada chamada. O poder está no \emph{flow}: o tipo de entrada
\emph{flui} para a saída, então \texttt{primeiroItem([10])} devolve
\texttt{number}, não \texttt{unknown}. Sem isso, você teria que lançar
\texttt{as number} — uma mentira para o compilador.
\end{caixaconceito}

\section{Tipos utilitários nativos}

O TypeScript traz utilitários prontos que derivam novos tipos de tipos
existentes — você os usará diariamente:

\begin{lstlisting}[style=livro, language=TypeScript,
  caption={Os utilitários de tipo mais usados.}, label={list:ts-utilitarios}]
interface Produto {
  id: number;
  nome: string;
  preco: number;
  descricao: string;
}

type SemDescricao = Omit<Produto, "descricao">;   // remove campos
type SomenteBasico = Pick<Produto, "id" | "nome">; // mantém campos
type Parcial = Partial<Produto>;                    // tudo opcional
type Imutavel = Readonly<Produto>;                  // tudo readonly
type Requerido = Required<Partial<Produto>>;        // tudo obrigatório

// Para APIs: o que o cliente envia vs o que o servidor devolve
type CriarProduto = Omit<Produto, "id">;   // POST body
type ProdutoExibido = Pick<Produto, "id" | "nome" | "preco">; // GET response

// Tipos condicionais úteis
type ElementoOuArray<T> = T extends unknown[] ? T[number] : T;
type N = ElementoOuArray<number[]>;   // number
type S = ElementoOuArray<string>;     // string
\end{lstlisting}

\section{Tipos para DOM e para APIs}

O TypeScript conhece o DOM: \texttt{document.querySelector} já retorna os
tipos certos (com \texttt{null} possível!), e \texttt{instanceof} faz o
narrowing:

\begin{lstlisting}[style=livro, language=TypeScript,
  caption={DOM tipado e cast seguro.}, label={list:ts-dom}]
// querySelector tipado via generic
const campo = document.querySelector<HTMLInputElement>("#email");
if (campo) {
  // aqui campo é HTMLInputElement, não null
  campo.value = campo.value.trim();
}

// Eventos tipados
formulario.addEventListener("submit", (evento) => {
  evento.preventDefault();
  const dados = new FormData(evento.currentTarget as HTMLFormElement);
});

// Para APIs externas, defina o contrato na borda
interface GithubUser {
  login: string;
  name: string | null;
  public_repos: number;
}

async function buscarUsuario(login: string): Promise<GithubUser> {
  const resposta = await fetch(`https://api.github.com/users/${login}`);
  if (!resposta.ok) throw new Error(`HTTP ${resposta.status}`);
  return (await resposta.json()) as GithubUser; // contrato na fronteira
}
\end{lstlisting}

\begin{caixaarmadilha}
O cast \texttt{as GithubUser} acima é um \emph{contrato assumido}: se a API
mudar o formato, o TypeScript não vai reclamar. Para dados que vêm de fora
(APIs, JSON, formulários), a validação \emph{em runtime} é obrigatória — você
usará Zod no capítulo~\ref{cap:orms} e nas Server Actions do
capítulo~\ref{cap:fullstack}. A regra: \emph{tipos na compilação, validação na
borda}.
\end{caixaarmadilha}

\section{O que \emph{não} fazer com \texttt{any}}

\begin{itemize}[leftmargin=*, itemsep=0.4em]
  \item \texttt{any} desliga a rede de segurança — o valor ``cai'' para JavaScript puro;
  \item \texttt{as any} em cadeia é a versão ``cala a boca, compilador'' — sempre perde informação;
  \item \texttt{unknown} é o substituto seguro: força você a fazer narrowing antes de usar;
  \item Em código legado, migre \texttt{any} $\rightarrow$ \texttt{unknown} $\rightarrow$ tipo real, progressivamente.
\end{itemize}

\begin{lstlisting}[style=livro, language=TypeScript,
  caption={any vs unknown: a diferença que salva.}, label={list:ts-any}]
function comAny(x: any) {
  return x.nome.toUpperCase(); // compila — e explode se x for null
}

function comUnknown(x: unknown) {
  // x.nome // Erro: property 'nome' does not exist on type 'unknown'
  if (typeof x === "object" && x !== null && "nome" in x) {
    const obj = x as { nome: string };
    return obj.nome.toUpperCase();
  }
  return "valor desconhecido";
}
\end{lstlisting}

% ============================================================================
\section{Projeto do capítulo: biblioteca de validação tipada}
\label{sec:projeto6}

\begin{caixaprojeto}
\textbf{Objetivo:} construir uma pequena biblioteca de validação de
formulários — \texttt{valida.ts} — que retorna erros tipados por campo, com
regras encadeáveis (obrigatório, e-mail, mínimo/máximo de caracteres, regex).

\textbf{Critérios de aceite:}
\begin{itemize}[leftmargin=*, itemsep=0.2em]
  \item \texttt{strict: true} e \texttt{tsc --noEmit} sem erros;
  \item API encadeável: \texttt{v.campo("nome").obrigatorio().minimo(3).maximo(80)};
  \item Resultado tipado: \texttt{\{ valido: boolean; erros: Record<string, string> \}};
  \item Generic no validador de regex e nos campos;
  \item Pelo menos 6 testes unitários — até o capítulo~\ref{cap:testes}, valide
  o validador com \texttt{node --test} (nativo, sem configuração); a partir de
  lá, portado para Vitest;
  \item Integração opcional: validar o formulário de contato do portfólio (capítulo~\ref{cap:universo}).
\end{itemize}
\end{caixaprojeto}

O esqueleto da biblioteca (Listagem~\ref{list:ts-validador}):

\begin{lstlisting}[style=livro, language=TypeScript,
  caption={Núcleo do validador encadeável.}, label={list:ts-validador}]
type Regra = (valor: string) => string | null; // null = passou

class Campo {
  private regras: Regra[] = [];

  obrigatorio(): this {
    this.regras.push((v) => (v.trim() ? null : "Campo obrigatório"));
    return this;
  }

  minimo(n: number): this {
    this.regras.push((v) =>
      v.trim().length >= n ? null : `Mínimo de ${n} caracteres`,
    );
    return this;
  }

  email(): this {
    const re = /^[^\s@]+@[^\s@]+\.[^\s@]+$/;
    this.regras.push((v) => (re.test(v.trim()) ? null : "E-mail inválido"));
    return this;
  }

  validar(valor: string): string | null {
    for (const regra of this.regras) {
      const erro = regra(valor);
      if (erro) return erro;
    }
    return null;
  }
}

export function campo(): Campo {
  return new Campo();
}

// Uso:
// const erro = campo().obrigatorio().minimo(3).maximo(80).validar("ab");
// console.log(erro); // "Mínimo de 3 caracteres"
\end{lstlisting}

\begin{caixadica}
O padrão \emph{fluent interface} (``encadeável'') acima é o mesmo usado por
Zod, Prisma e tantas bibliotecas profissionais: cada método retorna
\texttt{this} para permitir o encadeamento, e o tipo de retorno é preservado.
Você acabou de construir, na prática, um mini-Zod.
\end{caixadica}

\section{Exercícios de fixação}

\begin{exercicios}
  \item Qual a diferença entre JavaScript e TypeScript? O que acontece com os tipos na compilação?
  \item O que faz \texttt{strictNullChecks} e por que ele é importante?
  \item Escreva a interface\index{interface} de um \texttt{Produto} (id, nome, preco, emEstoque) e um array tipado dela.
  \item Crie um union type \texttt{FormaDePagamento} com 3 valores e uma função que o trate com narrowing.
  \item Escreva um generic \texttt{duplicar(lista)} que retorna a lista com os elementos em dobro.
  \item Explique a diferença entre \texttt{interface} e \texttt{type} e dê um caso de uso para cada.
\end{exercicios}

\section{Desafios}

\begin{desafios}
  \item \textbf{Utilitários na prática}: use \texttt{Pick}, \texttt{Omit}, \texttt{Partial} e \texttt{Readonly} para derivar 4 tipos novos de \texttt{Usuario} e explique cada um.
  \item \textbf{Discriminated union}: modele \texttt{Resposta} como
  \texttt{\{ tipo: "sucesso"; dados: T \} | \{ tipo: "erro"; mensagem: string \}}
  e escreva uma função que trate os dois casos com narrowing.
  \item \textbf{API tipada}: defina os tipos da resposta de uma API pública
  (ex.: GitHub) e faça um \texttt{fetch} com o resultado tipado e validado.
  \item \textbf{Validador completo}: estenda o \texttt{valida.ts} com
  \texttt{.regex(re)}, \texttt{.maximo(n)} e \texttt{.correspondente(outroCampo)}
  e escreva os testes.
\end{desafios}

\section{Exercícios de nível sênior}

\begin{seniores}
  \item \textbf{Tipos condicionais}: implemente
  \texttt{type ElementoOuArray<T> = T extends unknown[] ? T[number] : T;} e
  explique onde tipos condicionais aparecem em bibliotecas reais (pistas:
  \texttt{Parameters}, \texttt{ReturnType}).
  \item \textbf{Mapeamento de tipos}: crie, na mão, um tipo que transforma
  todas as propriedades de uma interface em \emph{opcionais} (como o
  \texttt{Partial} nativo), usando \texttt{keyof} e \texttt{in}.
  \item \textbf{Segurança na fronteira}: compare validação \emph{em tempo de
  compilação} (tipos) com validação \emph{em runtime} (Zod) e explique por que
  as duas são necessárias em uma API.
  \item \textbf{Inferência de strings}: crie o tipo
  \texttt{type Rotas = "/"| "/blog"| "/blog/[slug]"} e um helper tipado que
  monta URLs com parâmetros, garantindo que o slug seja passado.
\end{seniores}

\section{Armadilhas comuns}

\begin{itemize}[leftmargin=*, itemsep=0.4em]
  \item Usar \texttt{any} como ``fuga'' — ele desliga toda a rede de segurança;
  \item \texttt{as Tipo} agressivo (cast que mente para o compilador e explode em runtime);
  \item Ignorar \texttt{strict} em projetos novos (projetos legados sofrem disso);
  \item Anotar tudo em excesso (ruído) ou nada (imprudência) — anote as fronteiras;
  \item Confundir \texttt{interface} e \texttt{type} sem critério — escolha uma convenção.
\end{itemize}

\section{Resumo do capítulo}

\begin{itemize}[leftmargin=*, itemsep=0.3em]
  \item TypeScript = JavaScript + tipos estáticos, apagados na compilação;
  \item \texttt{strict: true} (com \texttt{strictNullChecks}) é não negociável;
  \item \texttt{interface}/\texttt{type}, unions, literals, narrowing e generics formam o núcleo;
  \item Discriminated unions modelam estados e resultados com segurança;
  \item Utilitários (\texttt{Pick}, \texttt{Omit}, \texttt{Partial}...) derivam tipos novos;
  \item Tipos na compilação, validação em runtime na borda (Zod no capítulo~\ref{cap:orms});
  \item Projeto entregue: \texttt{valida.ts}, uma biblioteca de validação encadeável e tipada.
\end{itemize}

\section{Referências oficiais para aprofundar}

\begin{itemize}[leftmargin=*, itemsep=0.3em]
  \item Documentação oficial do TypeScript: \url{https://www.typescriptlang.org/docs/}
  \item TypeScript Handbook (PT-BR): \url{https://www.typescriptlang.org/pt/docs/handbook/intro.html}
  \item Playground oficial: \url{https://www.typescriptlang.org/play}
  \item tsconfig completo referenciado: \url{https://www.typescriptlang.org/tsconfig}
\end{itemize}
